Decision-time intervention for autonomous experimentation agents is usually
studied as a prediction problem: estimate which checkpoints are likely to fail,
then act on the ranking.  We study the decision problem directly, at a single
online-reachable checkpoint where a controller must choose between continuing
the frozen agent policy and invoking an independent verifier whose result can
change the continuation.  Feasibility is observed before treatment. Verification cost is measured per
branch, while harm to runs that would otherwise succeed is a counterfactual
quantity estimated from paired outcomes, not an observed decision-time feature.  We specify a
prospective branch-data contract for this problem, specify qualification of the runtime boundary
that licenses causal reading of branch contrasts, and compare outcome-based
allocation (failure risk, utility deficit) with action-conditioned allocation
under matched information, capacity and branch budget.  A conservative
selection rule with an explicitly allocated one-sided bound is specified so
that the first study can report what protection costs in intervention rate and
value.  The paper contributes a research design, a measurement contract and an
engineering smoke path; no real-agent experimental result is reported here, and
every performance number is marked planned.  The ordinal-sufficiency theorem
that motivates the comparison is prior art and is cited as such; what is new is
the budgeted, feasibility- and harm-constrained allocation problem and the
matched-information test of fitted policies on it.
